\documentclass[11pt,a4paper]{article}
\usepackage{fontspec}
\setmainfont{Noto Serif}
\setsansfont{Noto Sans}
\usepackage{babel}
\babelprovide[main,import]{russian}
\usepackage{amsmath,amssymb,mathtools,booktabs,geometry,fancyhdr,microtype,xcolor,hyperref}
\geometry{left=21mm,right=21mm,top=21mm,bottom=21mm}
\definecolor{ink}{HTML}{173D57}
\hypersetup{colorlinks=true,urlcolor=ink,linkcolor=ink,pdftitle={BGS-F. Арифметические осцилляции при критическом кике}}
\pagestyle{fancy}\fancyhf{}\lhead{\small BGS-F}\rhead{\small Арифметические осцилляции при критическом кике}\cfoot{\thepage}
\setlength{\headheight}{15pt}\setlength{\parindent}{1em}\setlength{\parskip}{2.5pt}
\setlength{\emergencystretch}{3em}\tolerance=3000
\newcommand{\Tr}{\operatorname{Tr}}
\newcommand{\EE}{\mathbb E}
\newcommand{\R}{\mathbb R}
\newcommand{\Z}{\mathbb Z}
\newcommand{\C}{\mathbb C}
\newcommand{\si}{\mathrm{Sine}_2}
\newcommand{\eps}{\varepsilon}
\newcommand{\sinc}{\operatorname{sinc}}
\newcommand{\sect}[1]{\section*{\sffamily\color{ink}#1}}
\newcommand{\subt}[1]{\par\medskip\noindent\textbf{#1.}\quad}
\newcommand{\T}{\mathbb T}
\newcommand{\dist}{\operatorname{dist}}
\newcommand{\Ad}{\operatorname{Ad}}
\newcommand{\Zp}{\mathbb Z_7}
\newcommand{\Qp}{\mathbb Q_7}
\newcommand{\Q}{\mathbb Q}
\newcommand{\HS}{\mathrm{HS}}
\begin{document}
\noindent{\small\sffamily\color{ink} МАТЕМАТИЧЕСКАЯ ФИЗИКА\hfill 22 НОЯБРЯ 2026}
\begin{center}{\LARGE\sffamily Арифметические осцилляции\\[2pt]при критическом кике}\\[6pt]
{\large Связь соседних двойных узлов и спектр сжатия}
\end{center}

Для исходной модели BGS-F положим \(N=7^k\), \(m=3N/7\), \(U_N(t)=e^{itA_N}M_N\) и \(t=\kappa/\sqrt N\) при фиксированном \(\kappa\ne0\). Физическая сила тогда равна \(u_N=t/N=\kappa N^{-3/2}\). Линейные узлы \(M_N\) имеют фазы \(\lambda_j\) с шагом \(2\pi/m\) и проекторы \(E_j\). Через \(X\) обозначим параметр характера первого узла.

\subt{Точный межузловой блок}
При равномерном выборе \(j\) в классе \(X\equiv5\bmod7\) оба узла \(j,j+1\) двойные. Сингулярные числа блока
\[
  R_{N,j}=\frac{N}{i\lambda_j}E_j U_N(\kappa/\sqrt N)E_{j+1}
\]
не имеют общего предельного распределения. Вдоль любой подпоследовательности, на которой указанная ниже арифметическая фаза \(\eta_k\) стремится к \(\eta\), закон неупорядоченной пары равен
\[
  \frac{7|\kappa|}{600}\,\bigl|-2V+(1+\eta)W\bigr|\,
  \bigl(|\cos\Theta|,|\sin\Theta|\bigr).
\]
Здесь \(V,W\) независимы и равномерны на единичной окружности, а \(\Theta\) независим и равномерен на \([0,\pi]\). Осцилляция среднего квадрата нормы Гильберта--Шмидта составляет не меньше \(\kappa^2/5\,760\,000\), где осцилляция означает разность верхнего и нижнего пределов.

\subt{Эрмитово сжатие}
При равномерном выборе \(j\) в классе \(X\equiv5\bmod49\) рассмотрим четыре собственных значения сжатия эрмитова карандаша
\[
  2N\!\left[\frac{(M_N-\lambda_j)(M_N+\lambda_j)^{-1}}{i}
              +\tan(tA_N/2)\right]
\]
на \(\operatorname{ran}(E_j+E_{j+1})\). Их совместный закон тоже не имеет предела. Осцилляция среднего следа квадрата сжатия составляет не меньше \(\kappa^2/2\,880\,000\). Кроме того, хотя бы для одного из лагов \(1,\ldots,5\) средняя энергия межузлового блока осциллирует и при выборе \(j\) по всей угловой сетке.

\subt{Причина осцилляции}
Две вейлевы метки исходного символа имеют общую норму \(-97\). Их стационарные вклады складываются с множителем \(1+\eta_k\), который остаётся виден после усреднения. Фаза \(\eta_k\) строится из иррационального отношения двух \(7\)-адических логарифмов. Это отдельный эффект от закона, доказанного ранее для класса \(X\equiv3\bmod7\) с шагом семь. Спектр полного \(U_N\) при критической силе и гипотеза BGS-F при \(u=10^{-80}\) остаются открытыми.

\newpage
\sect{1. Где возникает межузловая связь}
Берём \(\rho^2=2\), \(\rho\equiv3\bmod7\), \(\gamma=3+2\rho\), \(g=\mathcal L(\gamma)/7\) и \(\alpha=-(3g)^{-1}\), где \(\mathcal L(x)=\log(x^6)/6\) на единицах \(\Zp\). Продолжаем \(\mathcal L\) мультипликативно-аддитивным правилом \(\mathcal L(7)=0\). Пишем \(e_n(y)=e^{2\pi i y/n}\) для целого остатка \(y\) по модулю \(n\). Метка \(v=(v_1,v_2)\) имеет расщеплённые координаты и норму
\[
  a=v_1+\rho v_2,\qquad b=(v_1-\rho v_2)/(-2\rho),\qquad
  q=ab=-Q(v)/(2\rho),\qquad Q(v)=v_1^2-2v_2^2 .
\]
Оператор метки действует как \(T(a,b)f(x)=e_N(b(x+a/2))f(x+a)\), а \(M_N\) с точностью до общей фазы действует как \(f(x)\mapsto f(\gamma^{-1}x)\). У \(\mathcal L(\gamma)\) точная оценка \(\nu_7=1\), поскольку \(\gamma\) имеет порядок \(3\) по модулю \(7\) и \(21\) по модулю \(49\). Получаем \(\alpha\equiv3\bmod7\).

Пусть \(\omega(3)=\zeta_3\) задаёт характер на \(\mathbb F_7^\times\). На каждой из двух орбит единиц, образованных квадратами и неквадратами по модулю 7, собственный вектор имеет координаты
\[
   m^{-1/2}\,\omega(x)^h e_N(X\mathcal L(x)),\qquad
   j\equiv h\pmod3,\qquad j\equiv -3gX\pmod{7^{k-1}},
\]
и ноль вне своей орбиты. Это следует подстановкой \(\gamma^{-1}\) в характер с учётом \(\omega(\gamma)=\zeta_3^2\). На слое \(7^d\Zp^\times\) порядок дилатации равен \(m/7^d\) вплоть до последнего слоя. Вклад слоёв вне единиц поэтому появляется лишь при \(7\mid j\), что эквивалентно \(7\mid X\). Сосед \(j+1\) имеет параметры \(X+\alpha,h+1\). В выбранном классе \(X\equiv5\), \(X+\alpha\equiv1\pmod7\), и описанные базисы полны.

Вклад \(T(a,b)\) в \(E_jA_NE_{j+1}\) имеет фазу
\[
 \Phi(x)=b(x+a/2)+(X+\alpha)\mathcal L(x+a)-X\mathcal L(x),\qquad
        b x(x+a)+\alpha x-Xa=0
\]
в критической точке. Для каждого \(x\in\{1,\ldots,6\}\) достаточно сначала проверить последнее равенство по модулю 7 и условие \(x+a\ne0\). Полная таблица результатов следующая. Орбиты отмечены знаками квадратности.
\begin{center}
\begin{tabular}{cccc}
\toprule
Строка & Столбец & \(v\) & \((x,x+a)\bmod7\)\\
\midrule
\(+\)&\(+\)& \(A=(1,0),\ B=(1,-7),\ C=(1,7)\) & \((1,2)\)\\
\(+\)&\(-\)& \(A,B,C\)& \((2,3)\)\\
\(-\)&\(+\)& \(-A,-B,-C\)& \((5,4)\)\\
\(-\)&\(-\)& \(-A,-B,-C\)& \((6,5)\)\\
\bottomrule
\end{tabular}
\end{center}
Для трёх положительных меток уравнение по модулю 7 сводится к \((x-1)(x-2)=0\) с обратимым множителем. Для двух меток, делящихся на 7, остаётся \(\alpha x=0\), и вклад исчезает. Полный перебор по коэффициентам исходного символа воспроизводится целочисленной программой. Все показанные точки просты. Разбиение суммы на остатки по \(7^{\lceil k/2\rceil}\) убивает классы без точки и даёт в каждом оставшемся классе фазу в точном подъёме корня, умноженную на \(\sqrt N\) и гауссов множитель модуля один. При чётном \(k\) множитель равен 1, при нечётном он равен \(7^{-1/2}\sum_{y\bmod7}e_7(\Phi''(x)y^2/2)\). Тамошний характер \(\omega(x+a)^{h+1}\overline{\omega(x)^h}\) постоянен на остаточном классе. Это следует из линейного разложения фазы на указанной сетке и квадратичного разложения на сетке \(7^{\lfloor k/2\rfloor}\) при нечётном \(k\).

\newpage
\sect{2. Точная интерференция двух меток}
Положим \(z=x/a\) для трёх положительных меток. Их стационарное уравнение и значение фазы принимают вид
\[
 \begin{split}
 qz(z+1)+\alpha z&=X,\\
 F_{q,a}(X)&=q(z+1/2)+(X+\alpha)\mathcal L(z+1)-X\mathcal L(z)
                  +\alpha\mathcal L(a).
 \end{split}
\]
При фиксированной ветви \(z\) фазы \(B,C\) различаются ровно на \(\alpha\mathcal L(G)\), где
\[
 G=\frac{1+7\rho}{1-7\rho},\qquad
 \eta_k=e_{7^k}\bigl(\alpha\mathcal L(G)\bigr)
       =e_{7^{k-1}}(\xi),\qquad
 \xi=-\frac{\mathcal L(G)}{3\mathcal L(\gamma)}\in\Zp .
\]
Две ветви \(A\) и \(B\) имеют одинаковые конечные гауссовы и характерные множители на соответствующих строках таблицы. Действительно, \(a,b,x,x+a\) у всех трёх меток совпадают по модулю 7. Замена \((v,x)\mapsto(-v,-x)\) оставляет каждый вклад фазы и конечный множитель тем же, поскольку \(\mathcal L(-x)=\mathcal L(x)\) и \(\omega(-1)=1\).

Сумма двух стационарных значений для одной метки выражается через произведения её корней. Имеем \(z_1z_2=-X/q\), \((1+z_1)(1+z_2)=-(X+\alpha)/q\) и \(q(z_1+z_2+1)=-\alpha\). Следовательно,
\[
 F_{q,a}^{(1)}+F_{q,a}^{(2)}
     =\Psi(X)+\alpha\bigl(\mathcal L(a/b)-1\bigr),\qquad
 \Psi(X)=(X+\alpha)\mathcal L(X+\alpha)-X\mathcal L(X).
\]
Отношения \(a/b\) для \(A,B\) равны \(-2\rho\) и \(-2\rho/G\). Обозначим через \(V,W\) нормированные полные вклады первой ветви \(A,B\) и через \(\Omega\) произведение вкладов обеих ветвей \(A\), включая конечные множители. Вторая ветвь \(A\) равна \(\Omega\bar V\), вторая ветвь \(B\) равна \(\Omega\eta_k^{-1}\bar W\). Каждую ветвь \(C\) получаем из соответствующей ветви \(B\) умножением на \(\eta_k\).

Коэффициенты трёх меток равны \(c_A=-2D,\ c_B=c_C=D=(6+8i)/4000\). В базисах из орбит квадратов и неквадратов точный нормированный блок \(B_{N,j}=\sqrt N\,E_jA_NE_{j+1}\) при \(k\ge3\) равен
\begin{equation}
 B_{N,j}=\frac7{1200}
   \begin{pmatrix} D/|D|&0\\0&\bar D/|D|\end{pmatrix}
   \begin{pmatrix} Z&\Omega\bar Z\\\Omega\bar Z&Z\end{pmatrix},
 \qquad Z=-2V+(1+\eta_k)W. \label{eq:block}
\end{equation}
Множитель \(7/1200\) получается из \(N/m=7/3\) и \(|D|=1/400\). Перестановка орбит во второй строке вносит сопряжение коэффициентов. Отдельно заметим \((1+\eta_k)\eta_k^{-1}=1+\bar\eta_k\), что объясняет общую форму обеих внедиагональных компонент.

Отсюда без асимптотики
\[
       \|B_{N,j}\|_{\HS}^2=4(7/1200)^2|Z|^2 .
\]
После удаления унитарного левого множителя оставшаяся матрица диагонализуется фиксированной симметрической и антисимметрической парой. Её сингулярные числа равны \((7/1200)|Z\pm\Omega\bar Z|\). Этот вывод инвариантен относительно выбора фаз собственных векторов и не требует согласования базисов между различными \(N\).

\newpage
\sect{3. Независимость трёх равномерных фаз}
На \(X\in5+7\Zp\) можно выбрать аналитические подъёмы каждой ветви по простому корню. Положим
\[
  b_q=-\frac{(q+\alpha)^2}{4q},\qquad
  Y_q^2=4q(X-b_q),\qquad
  z=\frac{-(q+\alpha)+Y_q}{2q}.
\]
Знак \(Y_q\) выбирается по рассматриваемой ветви. Дифференцирование стационарного значения даёт
\[
 F'_{q,a}=\mathcal L(z+1)-\mathcal L(z),\qquad
 F''_{q,a}=
 \frac12\left(\frac1{X+\alpha}-\frac1X\right)
 -\frac{2qX+\alpha(q+\alpha)}{2X(X+\alpha)Y_q},\qquad
 \Psi''=\frac1{X+\alpha}-\frac1X .
\]
Первое равенство использует стационарность по \(z\). Во втором можно взять \(z'=1/Y_q\) и привести \(-z'/[z(z+1)]\) к общему знаменателю. Для \(Q=1,-97\) точки \(b_q\) различны уже по модулю \(343\), соответственно \(143\) и \(290\). Каждая из функций \(Y_q^2\) имеет простой ноль только в своей точке и потому их квадратные классы независимы в \(\Qp(X)\). В порождённом ими поле существуют независимые смены знака радикалов. Если целочисленная комбинация \(n_1 F_{q_A,a_A}+n_2 F_{q_B,a_B}+n_3\Psi\) была бы постоянна на открытом диске, её вторая производная тождественно обнулилась бы в этом поле. Смены знака и ненулевые коэффициенты при \(Y_q^{-1}\) вынуждают \(n_1=n_2=0\). Затем \(\Psi''\ne0\) вынуждает \(n_3=0\). Переход от аналитического тождества к полевому допустим, поскольку подстановка ветвей в поле локальных аналитических зародышей инъективна.

\subt{Сокращение каждой ненулевой гармоники}
Фазы в переменной \(X\) зависят только от остатка по \(7^{k-1}\), поскольку при изменении этого остатка на кратное \(7^{k-1}\) исходный показатель умножает \(\mathcal L\) на это кратное, а \(\mathcal L\in7\Zp\). Можно усреднять по модулю \(7^k\) с той же нормировкой по условному классу. Пусть \(H\) есть нетривиальная комбинация из предыдущего абзаца. Для \(u\in7^s\Zp\), \(s\ge1\), выполнено
\[
       H(X+u)=H(X)+uH'(X)\pmod{7^{2s}} .
\]
Действительно, простой корень сдвигается на кратное \(u\). Разложение \(\mathcal L\) в окрестности единиц и стационарность по корню дают остаток указанной делимости. Для \(\Psi\) то же следует разложением в окрестности единиц \(X,X+\alpha\). Разбиваем \(X\) по модулю \(7^s\), \(s=\lceil k/2\rceil\), и суммируем по подъёмам. В каждой ячейке среднее \(e_N(H(X))\) равно нулю, если \(H'(X)\not\equiv0\bmod7^{k-s}\). Доля остальных ячеек стремится к нулю. В самом деле, производная аналитична и имеет лишь изолированные нули на компактном диске, поэтому её множество нулей имеет меру ноль. Конечные характерные и гауссовы множители постоянны по \(X\bmod7\) при каждом \(k,h\) и не мешают сокращению. Аргумент работает и на меньшем диске \(X\in5+49\Zp\).

По критерию Вейля \((V,W,\Omega)\) сходится к произведению трёх равномерных мер на окружности. Даже если \(\eta_k\) не имеет предела, разложение \(|Z|^2\) и сокращение перекрёстной гармоники дают
\begin{equation}
   \EE_j\|B_{N,j}\|^2_{\HS}
        =4(7/1200)^2(6+2\Re\eta_k)+o(1). \label{eq:energy}
\end{equation}
Вдоль \(\eta_k\to\eta\) независимый \(\Omega\) делает относительную фазу \(\Omega\bar Z/Z\) равномерной условно по \(Z\ne0\). Формула для сингулярных чисел распространяется непрерывно на \(Z=0\) и даёт заявленный закон.

\newpage
\sect{4. Иррациональность и нижняя граница осцилляции}
\subt{Иррациональность}
Если бы \(\xi\) была рациональна, то для некоторых целых \(a,b\), \(b>0\), выполнялось бы \(\mathcal L(G^b\gamma^{-a})=0\). Ядро \(\mathcal L\) на \(\Zp^\times\) состоит из шестых корней единицы, поскольку логарифм инъективен на \(1+7\Zp\). В вещественном поле \(\Q(\sqrt2)\) это только \(\pm1\). Получилось бы
\[
            (1+7\rho)^b= \pm\gamma^a(1-7\rho)^b.
\]
Обе стороны лежат в \(\Z[\rho]\), в том числе при отрицательном \(a\), так как \(\gamma^{-1}=3-2\rho\). В гомоморфизме \(\Z[\rho]\to\mathbb F_{97}\) с \(\rho\mapsto-1/7\) корень \(\rho^2=2\) корректен благодаря \(2\cdot49\equiv1\bmod97\). Левая сторона обнуляется, правая обратима. Это противоречие доказывает \(\xi\notin\Q\).

\subt{Действительная часть не сходится}
Пусть \(S\) есть множество предельных точек \(\eta_k\) на окружности. Соотношение \(\eta_{k+1}^7=\eta_k\) влечёт \(S^7=S\). Если бы \(\Re\eta_k\) сходилась, в \(S\) было бы не более двух точек. Возведение в седьмую степень переставляло бы их, так что их углы были бы рациональны со знаменателями, взаимно простыми с 7. Выбираем ближайшую к \(\eta_k\) точку \(w_k\in S\) для всех больших \(k\). Тогда \(w_k^7=w_{k-1}\) из отделённости точек, а угловые отклонения \(\varepsilon_k=\arg(\eta_k/w_k)/(2\pi)\to0\) удовлетворяют \(7\varepsilon_k=\varepsilon_{k-1}\) с некоторого места. Все \(\varepsilon_k\) рациональны. Поэтому \(J=7^{k-1}\varepsilon_k\) есть фиксированное рациональное число. Если \(r_k=[\xi]_{7^{k-1}}\), то \(r_k-J\) равно \(7^{k-1}\) умножить на рациональное число со знаменателем, взаимно простым с 7. При \(k\to\infty\) получаем \(\xi=J\) в \(\Qp\), что невозможно.

\subt{Количественная форма}
Множество предельных точек углов \(\arccos\Re\eta_k\) инвариантно и сюръективно под \(F(\theta)=\arccos(\cos 7\theta)\). Обозначим его диаметром \(d>0\), а через \(I\) обозначим его выпуклую оболочку в \([0,\pi]\). Если бы \(I\) не содержал внутреннего излома \(b=j\pi/7\), \(1\le j\le6\), то \(F\) на \(I\) растягивало бы диаметр ровно в семь раз, в противоречии с инвариантностью. Для такого \(b\in I\) имеем \(F(b)\in\{0,\pi\}\). Взяв предельную точку \(x\in I\), получаем
\[
  \frac{\pi}{7}\le |b-F(b)|\le |b-F(x)|+|F(x)-F(b)|\le d+7d .
\]
При этом осцилляция косинуса на концах \(I\) не меньше \(1-\cos d\), поскольку при фиксированном \(d\) минимум достигается у \(0\) или \(\pi\). Из \(1-\cos d\ge2d^2/\pi^2\) следует \(\operatorname{osc}_k\Re\eta_k\ge1/1568\). Подставив в \eqref{eq:energy}, получаем \(\operatorname{osc}_k\EE_j\|B_{N,j}\|^2_{\HS}\ge 1/5\,760\,000\).

\newpage
\sect{5. Перенос на точный унитарный оператор}
\subt{Средняя энергия одной вейлевой метки}
Пусть \(w\ne0\) есть фиксированная целочисленная метка. Её орбита под \(\Ad M_N\) имеет длину \(o_w\asymp_w N\). Действительно, хотя бы одна из двух расщеплённых координат имеет фиксированную конечную \(7\)-адическую оценку, а порядок \(\gamma\) по модулю \(7^r\) равен \(3\cdot7^{r-1}\). Вейлевы операторы этой орбиты ортогональны в норме Гильберта--Шмидта и имеют квадрат нормы \(N\). Проекция на подпространство с фазой \(e_m(-h)\) для \(\Ad M_N\) имеет вид
\[
  \frac1m\sum_{r=0}^{m-1}e_m(hr)(\Ad M_N)^rT(w)
                 =\sum_j E_jT(w)E_{j+h}.
\]
Квадрат её нормы равен \(N/o_w\) или нулю по условию согласования \(m\mid h o_w\). Это даёт \(\sum_j\|E_jT(w)E_{j+h}\|^2_{\HS}=O_w(1)\).

Оператор \(A_N^2\) есть конечная линейная комбинация таких меток с ограниченными коэффициентами. Это следует из формулы произведения Вейля, где появляются только суммы двух исходных меток и множители модуля один. Скалярная метка вообще не влияет на переходы \(h\ne0\bmod m\). Поэтому
\[
 \EE_j\|E_jA_N^2E_{j+h}\|_{\HS}^2=O(N^{-1})
\]
для равномерного выбора по всей сетке, а также на любом условном классе фиксированной положительной доли. Норма \(A_N\) ограничена суммой модулей коэффициентов символа.

\subt{Критический кик}
Для \(h=1\) разложение экспоненты в операторной норме даёт
\[
 R_{N,j}=
  e_m(1)\left(\kappa B_{N,j}
              +\frac{i\kappa^2}{2} E_jA_N^2E_{j+1}\right)+O(N^{-1/2})
\]
на паре двойных узлов. Остаток ограничен в обеих матричных нормах, поскольку размер блока фиксирован. Следовательно, \(R_{N,j}-\kappa e_m(1)B_{N,j}\) стремится к нулю в среднем квадрате. Собственные значения \((R_{N,j}^*R_{N,j})^{1/2}\) зависят липшицево от блока в операторной норме. Сингулярные законы из первого раздела и нижняя граница осцилляции энергии тем самым переносятся на точный \(U_N\). Нормы \(R_{N,j}\) равномерно ограничены на выбранном классе по \eqref{eq:block} и норме \(A_N\). Поэтому расхождение второго момента исключает слабую сходимость даже без предположений о равномерной интегрируемости.

\subt{Распознавание по размерностям}
Если выбирать \(j\) по всей сетке, событие \(X\equiv5\bmod7\) эквивалентно \(\dim E_{j+3}\ge4\). Действительно, у \(j+3\) параметр \(X+3\alpha\), и слои вне единиц появляются ровно при делимости на 7. Вероятность события равна \(1/7\). Поэтому совместный закон локальных блоков, включающий эту дискретную метку размерности, также не имеет предела. Для проверки достаточно взять ограниченную непрерывную функцию от сингулярных чисел, совпадающую на отмеченном классе с квадратом нормы Гильберта--Шмидта, и умножить её на индикатор события.

\newpage
\sect{6. Диагонали двухузлового сжатия}
Для контроля диагональных блоков \(H_r=\sqrt N\,E_{j+r}A_NE_{j+r}\), \(r=0,1\), сужаем начальный класс до \(X\equiv5\bmod49\). Здесь будем использовать \(L_r=X+r\alpha\). Остаются два полных орбитальных базиса. Для примитивного \(v\) стационарное уравнение в узле есть \(qz(z+1)=L_r\). Для \(r=0\) по модулю \(7\) вырождены только \(Q=1,-97\). У обеих меток соответствующая точка \(-q/4\equiv19\bmod49\). Поэтому дискриминант кратен \(7\) ровно один раз на выбранном классе. Стационарных корней по модулю \(49\) нет. Разбиение по \(7^{\lceil k/2\rceil}\) убивает эти вклады при \(k\ge4\). Для \(r=1\) активные примитивные корни просты. У меток \((0,\pm7)\) снимаем один множитель \(7\) из фазы. Приведённое уравнение для них даёт \(x^2\equiv2\) при \(r=0\) и \(x^2\equiv6\) при \(r=1\). Во втором случае вклада нет, в первом каждый простой корень даёт амплитуду \(\sqrt{7N}\) до нормировки.

Для каждой фиксированной нормы и ветви \(z\) стационарная фаза
\[
    F=q(z+1/2)+L_r\bigl(\mathcal L(z+1)-\mathcal L(z)\bigr)
\]
одинакова у всех меток этой группы. Внутри конкретного матричного элемента конечный характер тоже одинаков. Для примитивных меток это следует из \((x+a)/x=(z+1)/z\), для глубоких \((x+a)/x\equiv1\bmod7\). Квадратичные коэффициенты после перевода в одну координату \(z\) или \(7z\) отличаются квадратом единицы, поэтому гауссовы множители совпадают. Складываются коэффициенты меток данной нормы с тем же корнем и той же парой орбит. Противоположные ветви \(z\) и \(-1-z\) дают противоположные фазы. У каждой ненулевой диагональной фазы вторая производная равна
\[
    -\frac{1}{L_rY_q},\qquad Y_q^2=4q\bigl(X+q/4+r\alpha\bigr).
\]
Формула справедлива и при \(\nu_7(q)=2\), поскольку \(z=x/a\) может иметь отрицательную оценку, а логарифмическая производная сохраняет вид \(1/z\). Аналитичность и сокращение по методу раздела 3 можно проводить в исходной единичной координате \(x\) после снятия множителя \(7\).

Группы по нормам в каждом узле имеют разные точки ветвления. Между разными нормами перекрёстные гармоники исчезают по доказательству раздела 3. Для разных ветвей одной нормы разность фаз равна удвоенной фазе с точностью до знака, и её гармоника тоже исчезает. Отсюда
\[
 \EE_j\Tr H_r=o(1),\qquad
 \EE_j\Tr H_r^2=\frac{49}{9}\sum_{\text{группы }g}
  7^{d(g)}\left|\sum_{v\in g}c_v\right|^2+o(1).
\]
Здесь \(d(g)=1\) для глубоких меток и \(0\) для примитивных. Группировка делается отдельно для каждой пары орбит и каждой из двух ветвей. Коэффициенты берутся непосредственно из исходного символа.
\begin{center}
\begin{tabular}{cll}
\toprule
\(r\) & Активные нормы по возрастанию & Предел \(\EE_j\Tr H_r^2\)\\
\midrule
0 & \(-162,-127,-98,-71,-50,-8,-1\) & \(185222989/6000000\)\\
1 & \(-162,-127,-71,-50,-8,-1\) & \(2989/6000000\)\\
\bottomrule
\end{tabular}
\end{center}
Эти суммы вычисляются целыми коэффициентами в единицах \(1/4000\). В сертификате записаны также остатки точек ветвления по модулю \(2401\). Оценка амплитуд конечным числом гауссовых вкладов даёт \(\sup_{j,k}\|H_r\|<\infty\) на выбранном классе.

\newpage
\sect{7. Собственные значения эрмитова сжатия}
Пусть \(P=E_j+E_{j+1}\), \(z=\lambda_j\), \(t=\kappa/\sqrt N\), \(T=\tan(tA_N/2)\). При больших \(N\) тангенс определён, а \(M_N+zI\) обратима, так как \(m\) нечётно и \(z\) линейный узел. Тождество Кэли для \(U_N(t)=e^{itA_N}M_N\) имеет вид
\[
 \frac{(I-iT)(U_N(t)-zI)(M_N+zI)^{-1}}{i}
    = \frac{(M_N-zI)(M_N+zI)^{-1}}{i}+T.
\]
Правая часть эрмитова на единичной окружности. Без сжатия её сингулярность эквивалентна попаданию \(z\) в спектр \(U_N(t)\). При сжатии на \(P\) и умножении на \(2N\) получаем
\[
 \begin{split}
   K_{N,j}
   &=d_0 E_{j+1}+\kappa
       \begin{pmatrix} H_0&B_{N,j}\\B^*_{N,j}&H_1\end{pmatrix}
       +O(N^{-1/2}),\\
   d_0&=2\pi N/m=14\pi/3 .
 \end{split}
\]
Линейная часть на соседнем узле равна \(2N\tan(\pi/m)=d_0+O(N^{-2})\). Для тангенса \(2N\tan(tA_N/2)=NtA_N+O_\kappa(N^{-1/2})\) в операторной норме. При \(X\equiv5\bmod49\) все блоки равномерно ограничены.

Обозначим суммы из последней таблицы через \(c_0,c_1\). След квадрата раскрывается без смешанных межузловых членов первого порядка. Линейный член \(2d_0\kappa\EE_j\Tr H_1\) исчезает в пределе, и
\[
  \EE_j\Tr K_{N,j}^2=
  2d_0^2+\kappa^2\bigl(c_0+c_1+2\EE_j\|B_{N,j}\|_{\HS}^2\bigr)+o(1).
\]
Согласно \eqref{eq:energy}, осцилляция этого среднего не меньше \(\kappa^2/2\,880\,000\). При ограниченном спектре слабая сходимость закона четырёх собственных значений, упорядоченных по величине или рассматриваемых как мультимножество, вынуждала бы сходимость этого среднего. Получено противоречие.

\subt{Нули двухузлового карандаша}
Есть также интерпретация в терминах угловых корней самого сжатия. Заменим \(z\) на \(\lambda_j e^{ix/N}\). После сжатия продолжаем тангенсы через полюсы исключённых узлов. При фиксированном ограничении \(|x|\le L\) на двух выбранных узлах линейные тангенсы дают равномерно
\[
       K_{N,j}(x)=K_{N,j}(0)-xI_4+O_L(N^{-2}).
\]
Производная по \(x\) равна отрицательной диагональной матрице с элементами \(-\sec^2((\arg(\lambda_{j+r}/\lambda_j)-x/N)/2)\), \(r=0,1\). Для \(L\) больше равномерной нормы \(K_{N,j}(0)\) четыре упорядоченные собственные величины строго убывают от положительных до отрицательных значений, с кратностями считая возможные совпадения. Поэтому есть четыре угловых корня, каждый на расстоянии \(O_L(N^{-2})\) от соответствующего собственного значения \(K_{N,j}(0)\) в координате \(x\). Их закон также осциллирует. Эта часть использует только два выбранных узла. Нули полного карандаша дополнительно зависят от связи с его ортогональным дополнением и от дальних полюсов.

\newpage
\sect{8. Осцилляция без выбора исходного узла}
Рассмотрим полную энергию перехода \(W_N(h)=\sum_j\|E_jA_NE_{j+h}\|_{\HS}^2\). Орбитальная проекция из раздела 5 даёт её точно арифметически. Примитивные метки исходного символа группируются по \(Q\) и квадратности \(a\bmod7\). Действительно, такие орбиты определяются произведением \(ab\) и классом \(a\) по подгруппе \(\langle\gamma\rangle\), состоящей из единиц с квадратным остатком. Здесь \(Q\) различимы по модулю \(N\) при \(k\ge3\). В примитивной орбите длины \(m\) две метки на расстоянии \(n\) по орбите дают перекрёстный вес \(2\Re(c_v\bar c_w e_m(hn))\) после умножения \(W_N(h)\) на \(m/N\). Для \(Q=-1,-71,-127,-97\) суммы таких весов по обеим орбитам имеют соответственно вещественные амплитуды \(-112,-28,-28,400\) в единицах \(1/4000^2\). Сопряжённость и одинаковый шаг в противоположной орбите убирают мнимые части. Две метки глубины один идут по разным орбитам длины \(m/7\), поэтому вносят \(7\sum_{v=(0,\pm7)}|c_v|^2\) только при \(7\mid h\). Получаем
\[
 \frac mN W_N(h)=\frac{1233}{160000}+\frac{567}{200}\mathbf1_{7\mid h}
 +\frac{400\cos(h\theta_{97})-28\cos(h\theta_{71})
 -28\cos(h\theta_{127})-112\cos(2\pi h/m)}{16000000}.
\]
Углы \(\theta_D=2\pi n_D/m\) задаются равенствами \(\gamma^{n_D}=G_D\bmod N\), где
\[
     G_{97}=G,\qquad G_{71}=\frac{1+6\rho}{-1+6\rho},
     \qquad G_{127}=\frac{1+8\rho}{-1+8\rho}.
\]
Для \(Q=-1\) отношение равно \(\gamma\). Замена ориентации орбиты меняет знак угла и оставляет косинус. Поскольку \(G\equiv1\bmod7\), имеем \(3\mid n_{97}\), и \(n_{97}/3\equiv-\xi\bmod7^{k-1}\). Следовательно, \(\cos\theta_{97}=\Re\eta_k\).

Если бы \(W_N(h)\) сходилась для каждого \(h=1,\ldots,5\), то пределы существовали бы для первых пяти моментов знаковой меры
\[
  400\delta_{\cos\theta_{97}}-28\delta_{\cos\theta_{71}}
                                      -28\delta_{\cos\theta_{127}} .
\]
Действительно, \(\cos(h\theta)\) есть многочлен Чебышёва от \(\cos\theta\), а масса меры фиксирована. Разность любых двух подпоследовательных предельных мер имеет не больше шести точек носителя. Нулевые моменты степеней 0--5 зануляют каждый атом интерполяцией Лагранжа. Но у предельной меры ровно один положительный атом, даже при совпадениях носителей. Он определяет предел \(\cos\theta_{97}\), который по разделу 4 зависит от подпоследовательности. Значит для некоторого \(1\le h\le5\) нет сходимости усреднённой энергии.

Для тех же лагов средняя энергия точного блока \(N E_jU_N(\kappa/\sqrt N)E_{j+h}/(i\lambda_j)\) равна \(\kappa^2(N/m)W_N(h)+o(1)\). В разложении используется тот же среднеквадратичный расчёт \(A_N^2\), а остаток экспоненты в норме Гильберта--Шмидта ограничен через ранг два хотя бы одного из краевых проекторов. Это даёт осцилляцию среднего для точного оператора по всей сетке. Вывод об отсутствии слабого предела сделан выше для условного и отмеченного законов, где применим ограниченный тест.

\newpage
\sect{9. Проверки и граница результата}
\subt{Точное воспроизведение}
\texttt{critical\_oscillation/check.py} пересчитывает таблицу межузловых корней из исходных вейлевых меток, проверяет невырожденность и относительные фазы по модулям \(7^2,\ldots,7^6\). Там же проверяются рациональные коэффициенты в орбитальной формуле, вырожденные случаи по модулю \(49\), точки ветвления и точные средние диагональных энергий. \texttt{critical\_oscillation/data/certificates.json} содержит результаты и сравнивается побайтно в режиме \texttt{--verify}.

\subt{Независимое сравнение с матрицами}
\texttt{critical\_oscillation/check\_matrix.py} применяет исходные операторы сдвига и умножения непосредственно к орбитальным векторам. При \(N=343,2401\) максимальное отклонение в соседнем блоке от формулы \eqref{eq:block} меньше \(2\cdot10^{-15}\), а при \(N=2401\) отклонение диагональных формул меньше \(4\cdot10^{-15}\). При \(N=343\) дополнительно строятся точные конечномерные \(e^{itA_N}M_N\) и \(\tan(tA_N/2)\) численной диагонализацией, и проверяются явные границы остатка по \(t\). Орбитальная формула проверена независимым расчётом следов сопряжения дилатацией для лагов от 0 до 7. Эти численные проверки сверяют нормировки и знаки. Утверждения о пределах опираются на арифметическое доказательство.

\subt{Значение для критического режима}
Класс \(X\equiv5\bmod7\) имеет положительную плотность. Поэтому существование предела на разреженной угловой решётке \(X\equiv3\bmod7\), доказанное в следующей части отчёта, не переносится автоматически на совместный закон всех отмеченных блоков. Даже согласованный выбор фаз собственных векторов не устраняет найденную осцилляцию, поскольку наблюдаемые сингулярные числа и собственные значения инвариантны при таком выборе.

\subt{Оставшееся препятствие}
Убывающий кик \(t=\kappa/\sqrt N\) оставляет ненулевое фазовое перемешивание между близкими узлами в шкале \(1/N\). При расчёте настоящих фаз дополнение Шура добавляет к двухузловому карандашу матричную поправку, которая суммирует связи со всеми исключёнными узлами. Для вывода об осцилляции полного спектра надо контролировать эту поправку в том же масштабе, включая окрестности её полюсов. Полученные межузловые осцилляции сами по себе такого контроля не дают. При фиксированной физической силе \(10^{-80}\) параметр \(t=N10^{-80}\) растёт, поэтому разложение, использованное в настоящей части, не определяет этот предел.

\subt{Сверка узловых классов}
Основной вывод про сингулярные числа доказан на всём классе \(X\equiv5\bmod7\). Подкласс \(X\equiv5\bmod49\) используется для постоянной оценки диагонали в спектральном выводе. Сужение удаляет вырожденный подъём, а доказательство независимости фаз и нижняя граница осцилляции остаются теми же. Для глобальной средней энергии используется отдельная орбитальная формула по всей угловой сетке.

\end{document}
